\documentclass[11pt]{article}
\usepackage[margin=1in]{geometry}
\usepackage{amsmath,amssymb,amsthm}
\usepackage{booktabs}
\usepackage{xcolor}
\usepackage[colorlinks=true,linkcolor=blue,citecolor=blue,urlcolor=blue]{hyperref}

\newtheorem{theorem}{Theorem}[section]
\newtheorem{lemma}[theorem]{Lemma}
\newtheorem{proposition}[theorem]{Proposition}
\newtheorem{conjecture}[theorem]{Conjecture}
\newtheorem{question}[theorem]{Question}
\theoremstyle{remark}
\newtheorem{remark}[theorem]{Remark}

\newcommand{\todo}[1]{\textcolor{red}{\textbf{[TODO: #1]}}}
\newcommand{\Z}{\mathbb{Z}}
\newcommand{\F}{\mathbb{F}}
\newcommand{\cP}{\mathcal{P}}
\DeclareMathOperator{\lcm}{lcm}
\DeclareMathOperator{\ord}{ord}

\title{Toward a Counterexample to Agrawal's Conjecture\\
via Twin Smooth Integers\\
\large (working draft --- v0.1)}
\author{Dan Shumow \and Claude (Anthropic Fable~5)}
\date{August 6, 2026}

\begin{document}
\maketitle

\begin{abstract}
Agrawal's conjecture asserts that for coprime $n$ and prime $r$, the congruence
$(X-1)^n \equiv X^n-1 \pmod{n,\,X^r-1}$ forces $n$ prime or $n^2 \equiv 1
\pmod r$; if true it yields a deterministic $\tilde O(\log^3 n)$ primality
test. Lenstra and Pomerance gave a heuristic construction suggesting
counterexamples exist but are astronomically large; none is known, and
exhaustive search is hopeless. We pursue a \emph{constructive} search for a
counterexample of unrestricted size. Our contributions so far: (i) a reduction
of the Lenstra--Pomerance construction to a two-sided subset-product problem
over a partitioned smooth factor base; (ii) the observation that the candidate
pool consists exactly of \emph{twin smooth integers with prime sum} in a fixed
congruence class, connecting the search to the isogeny-cryptography parameter
literature and its datasets; (iii) feasibility estimates from that literature's
empirical supply curves, isolating a single unmeasured quantity --- the
``coloring tax'' --- as the go/no-go criterion; and (iv) an obstruction: the
boosting polynomials $2x^n-1$ of the twin-smooth literature can never satisfy
the local conditions of the Lenstra--Pomerance theorem as stated, motivating a
re-derivation of its admissible local conditions.
\end{abstract}

\section{Introduction}

For coprime integers $n, r$ write $T(-1,n,r)$ for the congruence
\begin{equation}\label{eq:aks}
(X-1)^n \equiv X^n - 1 \pmod{n,\, X^r-1}.
\end{equation}

\begin{conjecture}[Agrawal; Bhattacharjee--Pandey \cite{BP01}]\label{conj:agrawal}
If $r$ is a prime not dividing $n$ and $T(-1,n,r)$ holds, then $n$ is prime or
$n^2 \equiv 1 \pmod r$.
\end{conjecture}

If Conjecture~\ref{conj:agrawal} were true, a single evaluation of
\eqref{eq:aks} at one $r = O(\log n)$ would certify primality, giving a
deterministic $\tilde O(\log^3 n)$ test, versus $\tilde O(\log^6 n)$ for the
best proven algorithm \cite{LP2019}. Its status is strange: Lenstra and
Pomerance \cite{LPremarks} gave a heuristic argument for $\gg x^{1-\varepsilon}$
counterexamples below $x$, yet the conjecture has survived exhaustive search to
$10^{17}$ (Primaboinca, 2010--2020) and folklore places the least counterexample
at hundreds of digits. Popovych \cite{Popovych09} conjectured that adding the
hypothesis $(X+2)^n \equiv X^n+2 \pmod{n,\,X^r-1}$ repairs the statement; his
variant has no evidence against it.

\emph{Program.} We abandon minimality. Any counterexample refutes
Conjecture~\ref{conj:agrawal}; the L--P density heuristic \emph{improves} with
size; and the construction of \S\ref{sec:lp} is deterministic --- when its
arithmetic side conditions hold, \eqref{eq:aks} follows with no further luck
required. The size of a counterexample produced this way is bounded only by the
combinatorics of a subset-product problem (\S\ref{sec:subset}), whose raw
material turns out to be the \emph{twin smooth integers} of isogeny-based
cryptography (\S\ref{sec:twin}). Verifying a candidate $n$ of even $10^5$
digits directly against \eqref{eq:aks} costs a few core-hours, and any find
will also be tested against Popovych's second congruence.

\section{The Lenstra--Pomerance construction}\label{sec:lp}

\begin{theorem}[Lenstra--Pomerance \cite{LPremarks}; statement following
V\'a\v{n}a \cite{Vana09}, who proves the $k \equiv 3 \pmod 4$ case]
\label{thm:lp}
Let $n = p_1 \cdots p_k$ with $p_1, \dots, p_k$ distinct primes such that
\begin{enumerate}\itemsep1pt
\item[(a)] $k$ is odd;
\item[(b)] $p_i \equiv 3 \pmod{80}$ for all $i$;
\item[(c)] $p_i - 1 \mid n - 1$ for all $i$ \hfill ($n$ is Carmichael);
\item[(d)] $p_i + 1 \mid n + 1$ for all $i$ \hfill ($n$ is Lucas--Carmichael).
\end{enumerate}
Then $T(-1,n,5)$ holds and $n^2 \not\equiv 1 \pmod 5$; in particular such an
$n$ with $k \ge 3$ is a counterexample to Conjecture~\ref{conj:agrawal}.
\end{theorem}

\begin{proof}[Proof sketch]
Since $p \equiv 3 \pmod 5$ has order $4$ in $(\Z/5)^\times$, over $\F_p$ the
polynomial $X^5-1$ factors as $(X-1)\Phi_5(X)$ with $\Phi_5$ irreducible, so
$\Z[X]/(p, X^5-1) \cong \F_p \times \F_{p^4}$. The $\F_p$ component of
\eqref{eq:aks} is trivial. In $\F_{p^4}$, with $\zeta$ a primitive fifth root
of unity, the Frobenius gives $(\zeta-1)^{p^j} = \zeta^{p^j}-1$ for every $j$.
Conditions (c),(d) give $n \equiv 1 \pmod{p-1}$ and $n \equiv -1 \pmod{p+1}$,
i.e.\ $n \equiv p \equiv p^3 \pmod{\lcm(p-1,p+1)}$, and the mod-$80$ condition
resolves the $2$-adic ambiguity; V\'a\v{n}a's analysis of the quantity
$\rho(p) \mid 10(p^2-1)$ \cite[Thm.~3.4]{Vana09} shows control modulo
$10(p^2-1)$ suffices to conclude $(\zeta-1)^n = \zeta^n-1$ in each residue
field. Lemma~\ref{lem:mod5} gives $n^2 \not\equiv 1 \pmod 5$.
\todo{Phase 0: full independent re-derivation, both $k \bmod 4$ cases; see
Question~\ref{q:cells}.}
\end{proof}

\begin{lemma}\label{lem:mod5}
If $p_i \equiv 3 \pmod 5$ for all $i$ and $k$ is odd, then
$n \equiv 3^k \equiv \pm 2 \pmod 5$; hence $n^2 \equiv 4 \not\equiv 1 \pmod 5$.
\end{lemma}
\begin{proof}
$3^k \bmod 5$ cycles $3, 4, 2, 1$ with period $4$; odd $k$ gives $3$ or $2$,
i.e.\ $\pm 2 \pmod 5$.
\end{proof}

The difficulty of satisfying (c) and (d) simultaneously is concentrated in the
following trivial but decisive observation.

\begin{lemma}[Side-coprimality]\label{lem:sides}
If $n$ satisfies (c) and (d), then $\gcd(p_i-1,\,p_j+1) \mid 2$ for all $i,j$
(including $i=j$). Consequently the odd primes dividing
$\prod_i (p_i-1)$ and those dividing $\prod_i (p_i+1)$ form disjoint sets.
\end{lemma}
\begin{proof}
An odd prime $q$ dividing $p_i-1$ and $p_j+1$ divides both $n-1$ and $n+1$,
hence their difference $2$, a contradiction.
\end{proof}

\begin{remark}
No integer that is simultaneously Carmichael and Lucas--Carmichael is known;
none exists below $10^{18}$ \cite{Pinch}. Lemma~\ref{lem:sides} is the
structural reason, and the reason the \emph{least} counterexample is expected
to be enormous. It does not, however, obstruct \emph{large} solutions, which is
what we seek.
\end{remark}

\begin{lemma}[$2$-adic bookkeeping]\label{lem:2adic}
If $p \equiv 3 \pmod{16}$ then $v_2(p-1) = 1$ and $v_2(p+1) = 2$; if
$p \equiv 3 \pmod 5$ then $5 \nmid (p-1)(p+1)$. If moreover every $p_i \equiv 3
\pmod 4$ and $k$ is odd, then $n \equiv 3 \pmod 4$, so the $2$-power parts of
conditions (c),(d) hold automatically.
\end{lemma}
\begin{proof}
$p \equiv 3 \pmod{16}$ gives $p - 1 \equiv 2$ and $p+1 \equiv 4 \pmod{16}$,
fixing the stated $2$-adic valuations; $p \equiv 3 \pmod 5$ gives $p \mp 1
\equiv 2, 4 \pmod 5$. For the last claim, $n \equiv 3^k \equiv 3 \pmod 4$ for
odd $k$, so $2 \mid n-1$ with $v_2(n-1) \ge 1 = v_2(p_i-1)$ and
$4 \mid n+1$ with $v_2(n+1) \ge 2 = v_2(p_i+1)$.
\end{proof}

Thus under condition (b) the entire problem lives in odd moduli prime to $5$.

\section{Reduction to a two-sided subset-product problem}\label{sec:subset}

Fix disjoint finite sets $Q_-, Q_+$ of odd primes with $5 \notin Q_- \cup Q_+$,
and define the \emph{pool}
\[
\cP \;=\; \Bigl\{\, p \ \text{prime} : p \equiv 3 \!\!\pmod{80},\;
\tfrac{p-1}{2} \text{ is } Q_-\text{-smooth},\;
\tfrac{p+1}{4} \text{ is } Q_+\text{-smooth} \Bigr\}.
\]
Let $\Lambda_- = 2\,\lcm\{p-1 : p \in \cP\}$ and
$\Lambda_+ = 4\,\lcm\{p+1 : p \in \cP\}$.

\begin{proposition}\label{prop:reduction}
Let $S \subseteq \cP$ with $|S|$ odd, $|S| \ge 3$, and $n = \prod_{p \in S} p$.
If
\begin{equation}\label{eq:target}
n \equiv 1 \pmod{\Lambda_-} \qquad\text{and}\qquad n \equiv -1 \pmod{\Lambda_+},
\end{equation}
then $n$ satisfies conditions (a)--(d) of Theorem~\ref{thm:lp}, hence is a
counterexample to Conjecture~\ref{conj:agrawal}.
\end{proposition}
\begin{proof}
(a),(b) hold by construction of $\cP$ and $|S|$. For (c): $p_i - 1 \mid
\Lambda_-$ by definition of $\Lambda_-$, and $\Lambda_- \mid n-1$ by
\eqref{eq:target}. Likewise (d) via $\Lambda_+ \mid n+1$. The elements of $S$
are distinct primes, so $n$ is squarefree with $k = |S| \ge 3$, hence
composite. Note $\gcd(\Lambda_-, \Lambda_+) = 2$ by Lemma~\ref{lem:sides}
(disjointness of $Q_\pm$) and Lemma~\ref{lem:2adic}, so \eqref{eq:target} is
consistent.
\end{proof}

Condition \eqref{eq:target} is \emph{stronger} than (c),(d) --- it imposes the
full lcm rather than the lcm over $S$ --- but in exchange it is linear:
taking discrete logarithms in each cyclic component of
$(\Z/\Lambda_-)^\times \times (\Z/\Lambda_+)^\times$ (Pohlig--Hellman
\cite{PH78} is cheap, all component orders being smooth by design),
\eqref{eq:target} becomes a $0/1$ subset-sum over $\bigoplus_i \Z/m_i$, with
one extra $\Z/2$ coordinate enforcing $|S|$ odd. The heuristic solvability
condition is
\begin{equation}\label{eq:feas}
2^{|\cP|} \;\gg\; |G|, \qquad G = (\Z/\Lambda_-)^\times \times
(\Z/\Lambda_+)^\times,
\end{equation}
i.e.\ $|\cP| \gtrsim \log_2 \Lambda_- + \log_2 \Lambda_+$ with slack for the
solver. Candidate solvers, in escalating order: staged CRT in the style of
L\"oh--Niebuhr's Carmichael factories \cite{LN}, Wagner's $k$-list generalized
birthday algorithm \cite{Wagner} (designed for exactly this pool-rich regime),
and lattice reduction only for terminal low-dimensional corrections.
\todo{Phase 2: work out the staged-CRT schedule and the expected $|S|$.}

\section{The twin-smooth dictionary}\label{sec:twin}

\begin{lemma}[Dictionary]\label{lem:dict}
Let $m \ge 1$ and $p = 2m+1$. Then $p \in \cP$ (for $Q_\pm$ large enough to
contain the relevant primes) if and only if
\begin{enumerate}\itemsep1pt
\item[(i)] $m$ and $m+1$ are both smooth with respect to $Q_- \cup Q_+ \cup
\{2\}$, with the odd part of $m$ being $Q_-$-smooth and the odd part of $m+1$
being $Q_+$-smooth;
\item[(ii)] $m \equiv 1 \pmod{40}$;
\item[(iii)] $p$ is prime.
\end{enumerate}
Moreover (ii) alone implies $p \equiv 3 \pmod{80}$, $v_2(p-1)=1$,
$v_2(p+1)=2$, and $5 \nmid m(m+1)$.
\end{lemma}
\begin{proof}
$p - 1 = 2m$ and $p+1 = 2(m+1)$, so the smoothness conditions on
$(p-1)/2 = m$ and $(p+1)/4 = (m+1)/2$ are exactly (i) once the $2$-adic
valuations match. For (ii): $m \equiv 1 \pmod{40} \iff p = 2m+1 \equiv 3
\pmod{80}$. Then $m$ is odd (so $v_2(p-1)=1$), $m \equiv 1 \pmod 4$ gives
$m + 1 \equiv 2 \pmod 4$ (so $v_2(p+1) = 2$), and $m \equiv 1 \pmod 5$ gives
$m \not\equiv 0$, $m+1 \equiv 2 \not\equiv 0 \pmod 5$.
\end{proof}

Pairs of consecutive smooth integers $(m, m+1)$ --- \emph{twin smooths} ---
with $2m+1$ prime are precisely the parameters sought for the isogeny-based
schemes B-SIDH \cite{BSIDH} and SQIsign, and 2020--2025 produced an industrial
literature on generating them:

\begin{itemize}\itemsep2pt
\item \textbf{PTE sieving} (Costello--Meyer--Naehrig \cite{CMN21}):
Prouhet--Tarry--Escott solutions give degree-$n$ polynomial pairs differing by
a constant that split into linear factors; yielded $240$--$256$-bit primes $p$
with $p^2-1$ being $2^{15}$-smooth, up to $512$-bit at $2^{28}$-smooth.
\item \textbf{CHM} (Bruno et al.\ \cite{BCCEMNS23}, optimizing
Conrey--Holmstr\"om--McLaughlin \cite{CHM13}): iterative generation of
\emph{near-complete} sets of twin $B$-smooths. Headline data: $B = 547$ yields
$82{,}026{,}426$ pairs, the largest of $122$ bits; implementation public.
\item \textbf{Pell/St\o rmer} (Buzek--Hasan--Liu--Naehrig--Vigil \cite{BHLNV};
classical \cite{Stormer,Lehmer64}): for fixed $B$ the set of twin $B$-smooths
is \emph{finite}, computable via Pell equations $x^2 - 2Dy^2 = 1$; complete
sets known provably for $B \le 100$ and heuristically for $B \le 113$
\cite{MSW}.
\item \textbf{SVP records} (Mulder--Sterner--van Woerden \cite{MSW}): shortest
vectors in a prime-number lattice locate the \emph{largest} twin for a given
$B$ (records: $196$-bit at $B = 751$, $213$-bit at $B=997$), with an estimator
for the largest twin as a function of $B$.
\item \textbf{Boosting} (in \cite{BCCEMNS23,Sterner23}): evaluate
$p_n(x) = 2x^n - 1$ at twin-smooth arguments to reach cryptographic sizes.
\end{itemize}

By Lemma~\ref{lem:dict}, a $1/40$ congruence slice of \emph{any} twin-smooth
dataset, filtered for primality of $2m+1$, is a ready-made pool $\cP$ ---
before the one condition their literature does not consider: the global
side-partition of Lemma~\ref{lem:sides}.

\section{Feasibility estimates}\label{sec:feasibility}

\subsection{Supply and demand at $B = 547$}

\emph{Demand.} $\log_2 \Lambda_- + \log_2 \Lambda_+ \le \sum_{2 < q \le B}
\log_2 q^{e_q}$ where $q^{e_q}$ is the largest power appearing in any used
$p \mp 1$; the side-partition splits this sum but does not double it. With
$\pi(547) = 101$ primes and exponents capped by observed usage, this is
$\approx \theta(547)/\ln 2 \approx 750$ bits plus modest slack from powers of
small primes: call it $1000$--$1500$ bits, and less if the chosen subpool
touches only a sub-universe of primes.

\emph{Supply.} $82$M pairs at $B = 547$, $\div 40$ for the congruence slice
(Lemma~\ref{lem:dict}(ii)), $\div \ln p \approx 20$--$55$ for primality of
$2m+1$ at the dataset's typical sizes: roughly $5 \times 10^4$ pool candidates
\emph{from existing data}. \todo{Phase 1a: regenerate the dataset and replace
these estimates with exact histograms.}

The cushion between supply and demand is therefore a factor $\approx 25$--$40$.
Moreover it \emph{grows} with $B$: demand grows linearly
($\approx 1.44\,B$ bits) while empirical twin counts grow like $B^{\alpha}$
with $\alpha \approx 2.5$--$3$ (from $B = 113 \to 547$). St\o rmer finiteness
caps the supply at each fixed $B$, but the ceiling recedes superlinearly.

\subsection{The coloring tax}\label{sec:coloring}

The unmeasured factor is the side-partition. Each pool element consumes
$\omega_-(p) + \omega_+(p)$ prime-side assignments (the number of distinct odd
primes dividing $(p-1)/2$ and $(p+1)/4$), and all chosen elements must respect
one partition $Q_- \sqcup Q_+$. Under a uniformly random balanced partition, an
element survives with probability $\approx 2^{-(\omega_- + \omega_+)}$, so the
expected number of survivors is $\sum_{p \in \cP} 2^{-\omega(p)}$. At $B=547$,
bulk-of-distribution twins ($30$--$50$ bits) have $\omega_- + \omega_+ \approx
8$--$14$, predicting only \emph{tens} of survivors from $5 \times 10^4$
candidates --- two orders of magnitude short of the $\sim$$1500$--$3000$
needed.

The program therefore rests on the question:

\begin{question}\label{q:coloring}
By how much does an \emph{optimized} partition beat a random one? Survivors
correlate (elements sharing factor support survive together), supply is
thickest exactly where $\omega$ is smallest, and demand shrinks with the
effective prime universe. The complete St\o rmer sets ($B \le 113$) permit
computing the exact optimum at small scale; the CHM set at $B = 547$ permits
evaluating it at target scale. \todo{Phase 1a go/no-go: survival rate
$\gtrsim 5\%$ means the direct route is live; $\lesssim 0.1\%$ kills it at
fixed $B$.}
\end{question}

\subsection{Heuristic assumptions ledger}

\begin{center}
\begin{tabular}{@{}lll@{}}
\toprule
Statement & Status & Role \\
\midrule
$\gcd(p-1,p+1)=2$; Lemmas \ref{lem:mod5}--\ref{lem:dict} & theorems & structure \\
St\o rmer finiteness at fixed $B$ & theorem & supply ceiling \\
Local densities $P(q \mid p\mp1) = \tfrac1{q-1}$, exclusive & exact & constants only \\
Two-sided smoothness independence ($\rho(u_-)\rho(u_+)$) & conjectural & steering only \\
One-sided density $\#\{p : p-1 \text{ smooth}\} \sim \rho(u)\pi(x)$ &
open (Erd\H{o}s--Pomerance line) & steering only \\
Empirical supply curves (CHM, St\o rmer sets) & data & replaces both \\
\bottomrule
\end{tabular}
\end{center}

\section{An obstruction to boosting, and admissible local conditions}
\label{sec:boost}

The boosting polynomials $p_n(x) = 2x^n - 1$ of \cite{BCCEMNS23,Sterner23} are
structurally attractive here: for $n = 2$, $(p+1)/2 = x^2$ and $(p-1)/2 =
(x-1)(x+1)$ with $\gcd(x, x^2-1) = 1$, so the two sides have automatically
disjoint prime support --- Lemma~\ref{lem:sides} satisfied by construction ---
and the two-sided smoothness burden moves to $(x, x^2-1)$ at \emph{half} the
bit-size, where Dickman densities are far kinder. Unfortunately:

\begin{proposition}\label{prop:boost}
For every $n \ge 2$ and every integer $x$: $2x^n - 1 \not\equiv 3 \pmod{16}$.
In particular no value of any $p_n$, $n \ge 2$, satisfies condition (b) of
Theorem~\ref{thm:lp}. For $n = 2$ the obstruction is doubled: $x^2 \equiv 2
\pmod 5$ is also unsolvable, while for $n = 3$ the mod-$5$ condition is
satisfiable ($x \equiv 3 \pmod 5$) and only the $2$-adic obstruction remains.
\end{proposition}
\begin{proof}
$2x^n \equiv 4 \pmod{16}$ requires $x^n \equiv 2 \pmod 8$. If $x$ is even then
$4 \mid x^n$ (as $n \ge 2$), so $x^n \equiv 0$ or $4 \pmod 8$; if $x$ is odd
then $x^n$ is odd. In no case is $x^n \equiv 2 \pmod 8$. For the mod-$5$
claims: the squares modulo $5$ are $\{0,1,4\} \not\ni 2$, while $3^3 = 27
\equiv 2 \pmod 5$.
\end{proof}

The obstruction is to condition (b) \emph{as stated}, not (necessarily) to the
theorem's mechanism. Condition (b) fixes one convenient cell of local
conditions ($p \equiv 3 \bmod 16$ giving $v_2(p\mp1) = (1,2)$, and
$\ord_5(p) = 4$); V\'a\v{n}a's $\rho(p)$ framework suggests other cells may be
admissible, e.g.\ pools with a uniform but different $2$-adic pattern
$v_2(p+1) = v$ for $v > 2$, compensated in the mod-$2^k$ component of the
target congruence \eqref{eq:target}.

\begin{question}[Admissible local cells]\label{q:cells}
Determine the full set of local conditions (residues mod $2^k$ and mod $5$,
allowed $v_2$ patterns, and the required uniformity across the pool) under
which the conclusion of Theorem~\ref{thm:lp} survives. In particular: is any
admissible cell reachable by $p_n(x) = 2x^n - 1$ for some $n$, or by another
side-separating polynomial family? A positive answer would collapse the
coloring tax of \S\ref{sec:coloring} structurally and halve the effective
smoothness size.
\end{question}

\noindent For $n = 3$: $p = 2x^3-1$ has $v_2(p+1) = 1 + 3v_2(x)$, so $x \equiv
2 \pmod 4$ gives the uniform pattern $v_2(p+1) = 4$, $v_2(p-1) = 1$ ---
a natural first test case for Question~\ref{q:cells}. \todo{Phase 0.}

\section{Program}\label{sec:program}

\begin{enumerate}\itemsep1pt
\item[\textbf{P0.}] Independent re-derivation of Theorem~\ref{thm:lp};
resolution of Question~\ref{q:cells}; numerical unit tests of the congruence
machinery; prior-art due diligence (in particular whether the L--P construction
was ever run at scale).
\item[\textbf{P1a.}] Data-first go/no-go: regenerate the CHM $B=547$ set,
apply the dictionary filters, and answer Question~\ref{q:coloring} by direct
optimization; calibrate against the complete $B \le 113$ sets.
\item[\textbf{P1b.}] If needed, a native harvester (constructive enumeration
of one side, trial division on the other) at larger $B$ or targeted
congruences.
\item[\textbf{P2.}] Subset-product solver (staged CRT, Wagner), with the
parity coordinate.
\item[\textbf{P3.}] Assembly and dual independent verification of any find,
including the direct computation of $T(-1,n,5)$ and the Popovych congruence;
write-up (a quantified negative result at fixed $B$ is also publishable).
\end{enumerate}

\begin{thebibliography}{99}\itemsep1pt\small

\bibitem{AKS} M.~Agrawal, N.~Kayal, N.~Saxena, \emph{PRIMES is in P},
Ann.\ of Math.\ \textbf{160} (2004), 781--793.

\bibitem{BP01} R.~Bhattacharjee, P.~Pandey, \emph{Primality testing}, BTech
thesis, IIT Kanpur, 2001.

\bibitem{LPremarks} H.~W.~Lenstra, Jr., C.~Pomerance, \emph{Remarks on
Agrawal's conjecture}, AIM workshop notes, 2003.
\url{https://aimath.org/WWN/primesinp/articles/html/50a/}

\bibitem{Popovych09} R.~Popovych, \emph{A note on Agrawal conjecture},
Cryptology ePrint 2009/008.

\bibitem{Vana09} T.~V\'a\v{n}a, \emph{Agrawal's conjecture and Carmichael
numbers}, Comenius University Bratislava, 2009.

\bibitem{HD21} A.~Hegde, P.~Devaraj, \emph{Heuristics for the construction of
counterexamples to the Agrawal conjecture}, Springer PROMS \textbf{344}, 2021.

\bibitem{LP2019} H.~W.~Lenstra, Jr., C.~Pomerance, \emph{Primality testing
with Gaussian periods}, J.\ Eur.\ Math.\ Soc.\ \textbf{21} (2019), 1229--1269.

\bibitem{Pinch} R.~G.~E.~Pinch, \emph{The Carmichael numbers up to $10^{18}$},
arXiv:math/0604376.

\bibitem{AGP} W.~R.~Alford, A.~Granville, C.~Pomerance, \emph{There are
infinitely many Carmichael numbers}, Ann.\ of Math.\ \textbf{139} (1994),
703--722.

\bibitem{LN} G.~L\"oh, W.~Niebuhr, \emph{A new algorithm for constructing
large Carmichael numbers}, Math.\ Comp.\ \textbf{65} (1996), 823--836.

\bibitem{PH78} S.~Pohlig, M.~Hellman, \emph{An improved algorithm for
computing logarithms over $GF(p)$\dots}, IEEE Trans.\ Inf.\ Theory \textbf{24}
(1978), 106--110.

\bibitem{Wagner} D.~Wagner, \emph{A generalized birthday problem}, CRYPTO 2002.

\bibitem{BSIDH} C.~Costello, \emph{B-SIDH: supersingular isogeny
Diffie--Hellman using twisted torsion}, ASIACRYPT 2020.

\bibitem{CMN21} C.~Costello, M.~Meyer, M.~Naehrig, \emph{Sieving for twin
smooth integers with solutions to the Prouhet--Tarry--Escott problem},
EUROCRYPT 2021; ePrint 2020/1283.

\bibitem{BCCEMNS23} G.~Bruno, M.~Corte-Real Santos, C.~Costello,
J.~K.~Eriksen, M.~Meyer, M.~Naehrig, B.~Sterner, \emph{Cryptographic smooth
neighbors}, ASIACRYPT 2023; ePrint 2022/1439.

\bibitem{CHM13} J.~B.~Conrey, M.~A.~Holmstr\"om, T.~L.~McLaughlin,
\emph{Smooth neighbors}, Experiment.\ Math.\ \textbf{22} (2013), 195--202.

\bibitem{BHLNV} J.~Buzek, J.~Hasan, J.~Liu, M.~Naehrig, A.~Vigil,
\emph{Finding twin smooth integers by solving Pell equations},
arXiv:2211.04315.

\bibitem{Sterner23} B.~Sterner, \emph{Towards optimally small smoothness
bounds for cryptographic-sized twin smooth integers and their isogeny-based
applications}, SAC 2024; ePrint 2023/1576.

\bibitem{MSW} E.~Mulder, B.~Sterner, W.~van Woerden, \emph{Large smooth twins
from short lattice vectors}, ANTS XVII.

\bibitem{Stormer} C.~St\o rmer, \emph{Quelques th\'eor\`emes sur l'\'equation
de Pell $x^2 - Dy^2 = \pm 1$ et leurs applications}, Skr.\ Norske Vid.\ Akad.,
1897.

\bibitem{Lehmer64} D.~H.~Lehmer, \emph{On a problem of St\o rmer}, Illinois
J.\ Math.\ \textbf{8} (1964), 57--79.

\end{thebibliography}

\end{document}
